% GENERATED FILE. Edit canonical lessons, then run npm run generate:books.
% canonical-source-hash: fnv1a64:44ad03917118b20b
% canonical-lessons: PA-C346-kasur, PA-C346-lukauna, PA-C346-ilzam, PA-C346-majbur, PA-C346-afsos

\chapter{Owning Up}
\label{ch:pa-a2-owning-up}
\hlchaptermodality{346}

\begin{chapteropening}
\textbf{By the end of this chapter:} \emph{I can admit a fault and say I am sorry.}

Complete the last lesson of chapter 346: I can admit a fault and say I am sorry.
\end{chapteropening}

\section[\texorpdfstring{\pa{ਕਸੂਰ} (kasūr)}{ਕਸੂਰ (kasūr)}]{\pa{ਕਸੂਰ} (kasūr) — a fault; blame}

\label{lesson:PA-C346-kasur}



\begin{quote}
Before the new one: say the Punjabi for to settle a matter, then the Punjabi for a proposal.
\end{quote}

\subsection*{You'll want to know: \pa{ਕਸੂਰ}}
\textbf{\pa{ਕਸੂਰ}} — \emph{kasūr} — \textquotedblleft{}a fault; blame\textquotedblright{}.

\begin{quote}
\pa{ਇਹ} \pa{ਮੇਰਾ} \pa{ਕਸੂਰ} \pa{ਨਹੀਂ} \pa{ਹੈ।}
\end{quote}

\begin{quote}
\emph{eh merā kasūr nahīṁ hai.}
\end{quote}

\begin{quote}
This is not my fault.
\end{quote}

A masculine noun.

\begin{cousinweb}
\textbf{\pa{ਕਸੂਰ}} (\emph{kasūr}) is a loan, usually traced through Persian to Arabic \emph{quṣūr}, \textquotedblleft{}shortcoming\textquotedblright{}.
\end{cousinweb}

\begin{grammarlens}[title={What you've built}]
One more word for asking and answering.
\end{grammarlens}

\subsection*{Your turn}
\begin{itemize}
\raggedright
  \item \emph{Say it:} \emph{kasūr}
  \item \emph{Say it:} \emph{kasūr}, once more
  \item \emph{Say it:} \emph{eh merā kasūr nahīṁ hai.} — This is not my fault.
  \item \emph{Recall:} say the Punjabi for to settle a matter, then the Punjabi for a proposal, then say \emph{kasūr} again
\end{itemize}

\begin{tcolorbox}[breakable,colback=teal!4,colframe=teal!35!black,title={Before you move on}]
What does \pa{ਕਸੂਰ} mean? (\textquotedblleft{}A fault; blame\textquotedblright{}.) Say it once more.
\end{tcolorbox}
\section[\texorpdfstring{\pa{ਲੁਕਾਉਣਾ} (lukāuṇā)}{ਲੁਕਾਉਣਾ (lukāuṇā)}]{\pa{ਲੁਕਾਉਣਾ} (lukāuṇā) — to hide something; to keep something secret}

\label{lesson:PA-C346-lukauna}



\begin{quote}
Before the new one: say the Punjabi for a proposal, then the Punjabi for a fault.
\end{quote}

\subsection*{You'll want to know: \pa{ਲੁਕਾਉਣਾ}}
\textbf{\pa{ਲੁਕਾਉਣਾ}} — \emph{lukāuṇā} — \textquotedblleft{}to hide something; to keep something secret\textquotedblright{}.

\begin{quote}
\pa{ਮੈਂ} \pa{ਕੁਝ} \pa{ਨਹੀਂ} \pa{ਲੁਕਾਉਂਦਾ} \pa{ਹਾਂ।}
\end{quote}

\begin{quote}
\emph{maiṁ kujh nahīṁ lukāundā hāṁ.}
\end{quote}

\begin{quote}
I am not hiding anything.
\end{quote}

Past (I …): \pa{ਮੈਂ} \pa{ਚਾਬੀ} \pa{ਲੁਕਾਈ}, \emph{maiṁ chābī lukāī}.

\begin{cousinweb}
\textbf{\pa{ਲੁਕਾਉਣਾ}} (\emph{lukāuṇā}) is the causative partner of \emph{lukṇā}, \textquotedblleft{}to hide oneself\textquotedblright{}, which goes back to Sanskrit \emph{lupyate}, \textquotedblleft{}is lost, is removed\textquotedblright{}; CDIAL lists Punjabi \emph{lukṇā} under it (CDIAL 11083).
\end{cousinweb}

\begin{grammarlens}[title={What you've built}]
One more word for asking and answering.
\end{grammarlens}

\subsection*{Your turn}
\begin{itemize}
\raggedright
  \item \emph{Say it:} \emph{lukāuṇā}
  \item \emph{Say it:} \emph{lukāuṇā}, once more
  \item \emph{Say it:} \emph{maiṁ kujh nahīṁ lukāundā hāṁ.} — I am not hiding anything.
  \item \emph{Recall:} say the Punjabi for a proposal, then the Punjabi for a fault, then say \emph{lukāuṇā} again
  \item \emph{Write it:} copy \textbf{\pa{ਲੁਕਾਉਣਾ}} once, checking the whole word
\end{itemize}

\begin{tcolorbox}[breakable,colback=teal!4,colframe=teal!35!black,title={Before you move on}]
What does \pa{ਲੁਕਾਉਣਾ} mean? (\textquotedblleft{}To hide something; to keep something secret\textquotedblright{}.) Say it once more.
\end{tcolorbox}
\section[\texorpdfstring{\pa{ਇਲਜ਼ਾਮ} (ilzām)}{ਇਲਜ਼ਾਮ (ilzām)}]{\pa{ਇਲਜ਼ਾਮ} (ilzām) — an accusation; a charge}

\label{lesson:PA-C346-ilzam}



\begin{quote}
Before the new one: say the Punjabi for a fault, then the Punjabi for to hide something.
\end{quote}

\subsection*{You'll want to know: \pa{ਇਲਜ਼ਾਮ}}
\textbf{\pa{ਇਲਜ਼ਾਮ}} — \emph{ilzām} — \textquotedblleft{}an accusation; a charge\textquotedblright{}.

\begin{quote}
\pa{ਇਹ} \pa{ਇਲਜ਼ਾਮ} \pa{ਸੱਚਾ} \pa{ਨਹੀਂ} \pa{ਹੈ।}
\end{quote}

\begin{quote}
\emph{eh ilzām sachchā nahīṁ hai.}
\end{quote}

\begin{quote}
This accusation is not true.
\end{quote}

A masculine noun.

\begin{cousinweb}
\textbf{\pa{ਇਲਜ਼ਾਮ}} (\emph{ilzām}) is a loan, usually traced through Persian to Arabic \emph{ilzām}, \textquotedblleft{}binding someone\textquotedblright{}.
\end{cousinweb}

\begin{grammarlens}[title={What you've built}]
One more word for asking and answering.
\end{grammarlens}

\subsection*{Your turn}
\begin{itemize}
\raggedright
  \item \emph{Say it:} \emph{ilzām}
  \item \emph{Say it:} \emph{ilzām}, once more
  \item \emph{Say it:} \emph{eh ilzām sachchā nahīṁ hai.} — This accusation is not true.
  \item \emph{Recall:} say the Punjabi for a fault, then the Punjabi for to hide something, then say \emph{ilzām} again
  \item \emph{Write it:} copy \textbf{\pa{ਇਲਜ਼ਾਮ}} once, checking the whole word
\end{itemize}

\begin{tcolorbox}[breakable,colback=teal!4,colframe=teal!35!black,title={Before you move on}]
What does \pa{ਇਲਜ਼ਾਮ} mean? (\textquotedblleft{}An accusation; a charge\textquotedblright{}.) Say it once more.
\end{tcolorbox}
\section[\texorpdfstring{\pa{ਮਜਬੂਰ} (majbūr)}{ਮਜਬੂਰ (majbūr)}]{\pa{ਮਜਬੂਰ} (majbūr) — helpless; forced by circumstances}

\label{lesson:PA-C346-majbur}



\begin{quote}
Before the new one: say the Punjabi for to hide something, then the Punjabi for an accusation.
\end{quote}

\subsection*{You'll want to know: \pa{ਮਜਬੂਰ}}
\textbf{\pa{ਮਜਬੂਰ}} — \emph{majbūr} — \textquotedblleft{}helpless; forced by circumstances\textquotedblright{}.

\begin{quote}
\pa{ਮੈਂ} \pa{ਮਜਬੂਰ} \pa{ਹਾਂ} \pa{ਅਤੇ} \pa{ਅੱਜ} \pa{ਨਹੀਂ} \pa{ਆਵਾਂਗਾ।}
\end{quote}

\begin{quote}
\emph{maiṁ majbūr hāṁ ate ajj nahīṁ āvāṅgā.}
\end{quote}

\begin{quote}
I can't help it, and I won't come today.
\end{quote}

\begin{cousinweb}
\textbf{\pa{ਮਜਬੂਰ}} (\emph{majbūr}) is a loan, usually traced through Persian to Arabic \emph{majbūr}, \textquotedblleft{}compelled\textquotedblright{}.
\end{cousinweb}

\begin{grammarlens}[title={What you've built}]
One more word for asking and answering.
\end{grammarlens}

\subsection*{Your turn}
\begin{itemize}
\raggedright
  \item \emph{Say it:} \emph{majbūr}
  \item \emph{Say it:} \emph{majbūr}, once more
  \item \emph{Say it:} \emph{maiṁ majbūr hāṁ ate ajj nahīṁ āvāṅgā.} — I can't help it, and I won't come today.
  \item \emph{Recall:} say the Punjabi for to hide something, then the Punjabi for an accusation, then say \emph{majbūr} again
  \item \emph{Write it:} copy \textbf{\pa{ਮਜਬੂਰ}} once, checking the whole word
\end{itemize}

\begin{tcolorbox}[breakable,colback=teal!4,colframe=teal!35!black,title={Before you move on}]
What does \pa{ਮਜਬੂਰ} mean? (\textquotedblleft{}Helpless; forced by circumstances\textquotedblright{}.) Say it once more.
\end{tcolorbox}
\section[\texorpdfstring{\pa{ਅਫ਼ਸੋਸ} (afsos)}{ਅਫ਼ਸੋਸ (afsos)}]{\pa{ਅਫ਼ਸੋਸ} (afsos) — regret; it's a pity}

\label{lesson:PA-C346-afsos}



\begin{quote}
Before the new one: say the Punjabi for an accusation, then the Punjabi for helpless.
\end{quote}

\subsection*{You'll want to know: \pa{ਅਫ਼ਸੋਸ}}
\textbf{\pa{ਅਫ਼ਸੋਸ}} — \emph{afsos} — \textquotedblleft{}regret; it's a pity\textquotedblright{}.

\begin{quote}
\pa{ਮੈਨੂੰ} \pa{ਬਹੁਤ} \pa{ਅਫ਼ਸੋਸ} \pa{ਹੈ} \pa{ਕਿ} \pa{ਮੈਂ} \pa{ਨਹੀਂ} \pa{ਆਇਆ।}
\end{quote}

\begin{quote}
\emph{mainū̃ bahut afsos hai ki maiṁ nahīṁ āiā.}
\end{quote}

\begin{quote}
I am very sorry that I didn't come.
\end{quote}

A masculine noun.

\begin{cousinweb}
\textbf{\pa{ਅਫ਼ਸੋਸ}} (\emph{afsos}) is a loan, usually traced to Persian \emph{afsōs}, \textquotedblleft{}regret\textquotedblright{}.
\end{cousinweb}

\begin{grammarlens}[title={What you've built}]
One more word for asking and answering.
\end{grammarlens}

\subsection*{Your turn}
\begin{itemize}
\raggedright
  \item \emph{Say it:} \emph{afsos}
  \item \emph{Say it:} \emph{afsos}, once more
  \item \emph{Say it:} \emph{mainū̃ bahut afsos hai ki maiṁ nahīṁ āiā.} — I am very sorry that I didn't come.
  \item \emph{Recall:} say the Punjabi for an accusation, then the Punjabi for helpless, then say \emph{afsos} again
  \item \emph{Write it:} copy \textbf{\pa{ਅਫ਼ਸੋਸ}} once, checking the whole word
\end{itemize}

\begin{tcolorbox}[breakable,colback=teal!4,colframe=teal!35!black,title={Before you move on}]
What does \pa{ਅਫ਼ਸੋਸ} mean? (\textquotedblleft{}Regret; it's a pity\textquotedblright{}.) Say it once more.
\end{tcolorbox}
